\chapter{Jesus Cristo foi Concebido pelo Poder do Espírito Santo, nasceu da Virgem Maria}
\chaptermark{Jesus Cristo nasceu da Virgem Maria}

O terceiro artigo da Profissão de Fé concentra-se no mistério central da fé cristã: a Encarnação do Filho de Deus. ``Jesus Cristo foi concebido pelo poder do Espírito Santo, nasceu da Virgem Maria'' resume o acontecimento pelo qual o Verbo eterno assumiu a natureza humana para nossa salvação. Este artigo revela não apenas o facto histórico do nascimento de Jesus, mas o mistério da união hipostática e o papel singular da Virgem Maria no plano da salvação. Os parágrafos 456 a 570 do Catecismo da Igreja Católica oferecem a base para compreender a profundidade deste mistério, unindo a doutrina da Encarnação, os dogmas marianos e a contemplação de toda a vida de Cristo como mistério salvífico.

\section{O Filho de Deus se fez homem (456--483)}

O Filho de Deus fez-se homem para nos salvar, reconciliando-nos com o Pai. A Encarnação não foi um acontecimento acidental, mas o cumprimento do plano eterno de amor de Deus. O Verbo eterno, consubstancial ao Pai, assumiu a natureza humana para curar a nossa natureza ferida pelo pecado. Esta verdade fundamental ilumina todo o mistério cristão.

A Encarnação revela o amor de Deus de forma suprema. ``Deus amou de tal modo o mundo que lhe deu o seu Filho único''. O Filho de Deus não assumiu uma humanidade abstracta, mas uma humanidade concreta, semelhante à nossa em tudo excepto no pecado. Deste modo, Ele se tornou o modelo perfeito de santidade para a humanidade.

A união hipostática é o dogma central que explica este mistério. Na pessoa única do Verbo divino subsistem duas naturezas --- divina e humana --- sem confusão, sem mudança, sem divisão e sem separação. Esta doutrina foi definida contra as heresias que negavam a verdadeira humanidade ou a verdadeira divindade de Cristo. A Igreja professa que Jesus Cristo é verdadeiro Deus e verdadeiro homem.

A Encarnação permite-nos conhecer o amor de Deus de modo visível e concreto. O Filho de Deus feito homem é o caminho para o Pai. Nele, a humanidade encontra o seu modelo de vida e a possibilidade de se tornar filho adoptivo de Deus. Este mistério continua a ser fonte de admiração e gratidão para os cristãos de todos os tempos.

O Filho de Deus feito homem transforma a compreensão da dignidade humana. Ao assumir a nossa carne, Deus elevou a condição humana a uma dignidade insuspeitada. O corpo humano, a vida familiar, o trabalho e o sofrimento adquirem novo significado à luz da Encarnação.

\section{``\ldots Concebido pelo poder do Espírito Santo, nascido da Virgem Maria'' (484--511)}

A conceição de Jesus pelo poder do Espírito Santo e o seu nascimento da Virgem Maria são o início visível do mistério da Encarnação. O Espírito Santo fecundou o seio virginal de Maria, fazendo com que ela concebesse o Filho eterno do Pai. Este acontecimento marca o início da ``plenitude dos tempos'' e realiza as promessas do Antigo Testamento.

A Virgem Maria desempenha um papel único e insubstituível no plano da salvação. Segundo a Constituição Dogmática \emph{Lumen Gentium}, Maria, como Mãe de Deus e Mãe da Igreja, coopera de modo singular na obra da redenção. Ela não é apenas o instrumento passivo da Encarnação, mas a serva do Senhor que, com o seu ``fiat'', consente livremente no plano divino. A Igreja vê nela o modelo perfeito de fé e de obediência.

O dogma da união hipostática é inseparável deste mistério. Uma só pessoa divina, a segunda Pessoa da Santíssima Trindade, subsiste em duas naturezas completas. Esta verdade foi defendida contra o nestorianismo, que separava as duas naturezas, e contra o monofisismo, que as confundia numa só. A Igreja professa que Maria é verdadeiramente Mãe de Deus (Theotokos), pois deu à luz a pessoa divina do Verbo encarnado.

Os quatro dogmas marianos expressam a grandeza da vocação de Maria: a maternidade divina, a virgindade perpétua, a Imaculada Conceição e a Assunção ao Céu. Cada um destes dogmas protege e ilumina o mistério da Encarnação. As heresias contrárias, como o arianismo ou o nestorianismo, diminuíam ou a divindade de Cristo ou a dignidade de Maria como Mãe de Deus.

O nascimento da Virgem Maria revela a humildade de Deus. O Criador do universo nasce numa gruta, envolto em faixas, colocado numa manjedoura. Este paradoxo do Deus-menino convida-nos a contemplar o amor que se faz pequeno para nos elevar. Maria, como primeira discípula, guarda e medita todos estes acontecimentos no seu coração.

\section{Os mistérios da vida de Cristo (512--570)}

Os mistérios da vida de Cristo não são meros factos do passado, mas realidades salvíficas que continuam a actuar na vida da Igreja e dos fiéis. Toda a vida de Jesus é mistério de redenção. A Igreja convida os cristãos a contemplar cada momento da existência de Cristo, desde a sua conceição até à sua morte e ressurreição, para participar dos seus méritos e se conformar cada vez mais com Ele.

A vida oculta de Jesus em Nazaré constitui o primeiro grande mistério. Durante trinta anos, o Filho de Deus viveu na simplicidade de uma vida familiar, partilhando o trabalho, a obediência e o crescimento progressivo. Esta longa ocultação revela o valor do quotidiano, do trabalho e da vida familiar. Jesus santificou a vida ordinária, mostrando que a santidade é possível em todas as circunstâncias.

O baptismo de Jesus no Jordão marca o início da sua vida pública. O Espírito Santo desce sobre Ele e o Pai manifesta o seu agrado. Jesus assume o seu papel de Servo sofredor e começa a revelar o Reino de Deus. O seu baptismo prefigura o nosso baptismo, no qual somos também ungidos pelo Espírito e declarados filhos amados do Pai.

A tentação no deserto revela a solidariedade de Jesus com a humanidade pecadora. Ele vence as tentações que Adão não soube vencer, mostrando que a obediência ao Pai é o caminho da vitória sobre o mal. Os quarenta dias no deserto preparam o seu ministério e são modelo para a vida quaresmal da Igreja.

Durante o seu ministério público, Jesus anuncia o Reino de Deus com palavras e obras. Os seus ensinamentos, especialmente as Bem-aventuranças e as parábolas, revelam o rosto misericordioso do Pai. Os milagres e os sinais, particularmente as curas e as expulsões de demónios, manifestam que o Reino de Deus chegou. Jesus associa os Apóstolos à sua missão e forma o núcleo da futura Igreja.

A Transfiguração no monte Tabor revela a glória divina de Jesus aos seus discípulos. Neste mistério, a humanidade de Cristo resplandece com a glória da sua divindade. A voz do Pai repete a ordem: ``Escutai-O!''. A Transfiguração fortalece a fé dos discípulos diante da Paixão que se aproxima e revela o destino final da humanidade redimida.

Os mistérios da Paixão e da Morte de Jesus culminam a sua missão. A agonia no Getsémani, o julgamento injusto, a flagelação, a coroação de espinhos e a crucificação revelam até que ponto chegou o amor de Deus. Jesus assume voluntariamente o sofrimento por amor à humanidade. A sua morte na cruz é o sacrifício que reconcilia o mundo com Deus.

A Ressurreição de Cristo é o mistério culminante que ilumina todos os anteriores. A vitória sobre a morte confirma tudo o que Jesus ensinou e fez. A sua ressurreição gloriosa é a garantia da nossa própria ressurreição. Através dos mistérios da vida de Cristo, o cristão é chamado a morrer para o pecado e a ressuscitar para uma vida nova.

A contemplação dos mistérios da vida de Cristo não é um exercício meramente intelectual. É uma participação real na vida do Senhor. A liturgia, especialmente o ciclo do Natal e da Páscoa, actualiza estes mistérios para que os fiéis possam viver deles. A devoção ao Santo Rosário é um meio privilegiado para meditar estes mistérios e conformar a própria vida à de Cristo.

A vida inteira de Jesus é uma manifestação progressiva do mistério pascal. Desde o primeiro instante da Encarnação até à Ascensão, tudo converge para a glorificação do Pai e a salvação dos homens. A Igreja, ao celebrar estes mistérios, torna presente a obra salvífica de Cristo para todos os tempos e lugares.

\section{A União Hipostática e o Papel de Maria}

A doutrina da união hipostática ilumina todo o terceiro artigo do Credo. Uma só pessoa divina subsiste em duas naturezas perfeitas. Esta verdade, definida nos Concílios de Éfeso e Calcedónia, protege o mistério da Encarnação contra distorções. Maria, ao dar à luz Jesus, é verdadeira Mãe de Deus, pois a pessoa nascida dela é a segunda Pessoa da Santíssima Trindade.

\emph{Lumen Gentium}, no capítulo VIII, apresenta Maria como membro eminentíssimo da Igreja e tipo da Igreja. Ela precedeu os cristãos no caminho da fé, da obediência e da união com Cristo. O seu ``sim'' na Anunciação é o modelo do consentimento que cada baptizado deve dar ao plano de Deus.

Os quatro dogmas marianos encontram aqui a sua unidade profunda. A Imaculada Conceição preparou Maria para ser digna Mãe do Salvador. A virgindade perpétua manifesta a novidade radical da Encarnação. A maternidade divina e a Assunção gloriosa confirmam a sua associação singular à obra redentora de Cristo.